\section{绪论}

\subsection{研究背景与意义}

【阐述研究背景与选题意义，建议 2～3 段。】

随着……技术的快速发展，……在……领域得到了广泛应用\upcite{lecun2015deep}。然而，现有方法仍存在以下不足：其一，……；其二，……；其三，……。因此，针对上述问题开展研究，具有重要的理论意义与实际应用价值。

\subsection{国内外研究现状}

【梳理国内外相关工作，简述代表性方法，并在末尾指出其不足。】

针对该问题，文献\upcite{vaswani2017attention}提出了……方法，通过……实现了……；文献\upcite{devlin2019bert}进一步……。国内学者在……等方面也开展了大量研究\upcite{zhongwen2020shili}。综合来看，现有研究在……方面仍存在改进空间，这正是本文工作的切入点。

\subsection{本文主要工作}

本文围绕【研究主题】展开研究，主要工作包括以下三个方面：

\begin{enumerate}
  \item 针对【问题一】，提出了【方法或模型名称】，通过【关键技术】解决【具体困难】；
  \item 针对【问题二】，设计了【模块或策略名称】，实现了【预期效果】；
  \item 在【数据集或实验平台】上完成了对比实验与消融实验，验证了所提方法的有效性。
\end{enumerate}

\subsection{论文组织结构}

全文共分为五章：第 1 章绪论，介绍研究背景与意义，梳理国内外研究现状，并给出本文的主要工作与组织结构；第 2 章相关理论与技术基础，介绍本文涉及的基础理论与关键技术；第 3 章方法设计，给出总体设计方案与各关键模块的具体设计；第 4 章实验与结果分析，介绍实验设置，并通过对比实验与消融实验验证所提方法的有效性；第 5 章总结与展望，总结全文工作并指出未来研究方向。
